\documentclass[conference]{IEEEtran}
\usepackage{cite}
\usepackage{amsmath,amssymb,amsfonts}
\usepackage{graphicx}
\usepackage{textcomp}
\usepackage{booktabs}
\usepackage{float}

\begin{document}

\title{MakFleet: A Semantic-Aware Spatio-Temporal Knowledge Graph
for Bodaboda Safety at Makerere University}

\author{
\IEEEauthorblockN{Ddumba Jonah, Ssebwana John Paul,
Sebunya Ronaldo, and Namuli Evelyn}
\IEEEauthorblockA{Department of Computer Science, College of Computing
and Information Sciences\\
Makerere University, Kampala, Uganda\\
\{jonah.ddumba, johnpaul.ssebwana, ronaldo.sebunya,
evelyn.namuli\}@students.mak.ac.ug}
}

\maketitle

\begin{abstract}
Conventional relational data warehouses are fundamentally ill-suited
to the high-frequency, erratic spatio-temporal dynamics of motorcycle
taxis known as bodabodas on university campuses. This paper presents
MakFleet, a Semantic-Aware Spatio-Temporal Knowledge Graph deployed
for the bodaboda network at Makerere University, Kampala, Uganda.
The system ingests IoT telematics at 5-second intervals, applies
Hidden Markov Model map-matching to snap noisy GPS coordinates to
the real campus road network, and enriches raw sensor readings into
semantically meaningful driving events. These events are stored in an
ontology-driven Neo4j Knowledge Graph comprising 212 intersections,
476 road segments, and 826,260 telemetry records across two academic
semesters. A custom Spatio-Temporal Graph Neural Network combining
Graph Convolutional Network spatial encoding and Long Short-Term
Memory temporal aggregation achieves an AUC of 0.9268 on validation
and 0.9471 on a temporal out-of-distribution test, with perfect
recall across all evaluation splits. A causal Business Intelligence
dashboard answers why a driver was flagged rather than merely what
was recorded. The system is publicly deployed and the dataset is
openly available.
\end{abstract}

\begin{IEEEkeywords}
Bodaboda Telemetry, Knowledge Graph, Spatio-Temporal GNN,
Semantic Data Engineering, Anomaly Detection, Campus Safety
\end{IEEEkeywords}


\section{Introduction}

Bodabodas are the primary mode of short-distance intra-campus
transport at Makerere University, Kampala. They carry students
between halls, libraries, and lecture rooms rapidly and cheaply,
but operate with no formal safety monitoring. Reckless driving
incidents---harsh braking, speeding, rapid acceleration near
pedestrian zones---go entirely undetected and unrecorded.

Traditional data warehousing approaches based on the Kimball Star
Schema are structurally inadequate for this problem. Relational
tables cannot encode the graph-like topology of campus road
networks, and expensive SQL joins create prohibitive latency for
real-time safety interventions \cite{kimball2013}. Furthermore, raw
sensor readings carry no semantic meaning: a sudden deceleration
near a zebra crossing is a safe stop, while the same signal on an
open road is reckless braking. Resolving this distinction requires
a spatially-aware ontological layer---not a tabular classifier.

This paper presents MakFleet, which reconceptualises the campus
fleet data warehouse as an epistemic, semantic infrastructure.
The key contributions are: (1) a three-phase semantic data
engineering pipeline with cryptographic provenance tracking; (2) an
ontology-driven Spatio-Temporal Knowledge Graph with a novel
ContextualStop node label that enforces the safe-stop/reckless-brake
distinction at the storage layer; (3) a custom GCN+LSTM
Spatio-Temporal Graph Neural Network (ST-GNN) trained under strict
semester-based splits; and (4) a causal Business Intelligence
dashboard that generates subgraph evidence for flagged anomalies.


\section{Related Work}

Spatio-temporal data warehouses have been proposed to extend
relational schemas with spatial data types, but fundamental
scalability limits emerge when graph topology must be co-modelled
alongside time series \cite{pelekis2014}. Graph Neural Networks
outperform tabular models for transport anomaly detection because
they encode road network topology directly \cite{zhu2022}. Yu
et al.'s STGCN \cite{yu2018} and Zhao et al.'s T-GCN \cite{zhao2019}
demonstrated strong spatio-temporal forecasting but assume static
city-scale road graphs and do not address semantic event enrichment.
Privacy vulnerabilities in GPS trajectory datasets are well
established; spatial k-anonymity and differential privacy are the
primary mitigations \cite{demontjoye2013}. Table~\ref{tab:related}
summarises key related works and their limitations relative to
MakFleet.

\begin{table}[t]
\caption{Summary of Related Works}
\label{tab:related}
\centering
\small
\begin{tabular}{|p{1.6cm}|p{1.4cm}|p{1.7cm}|p{1.9cm}|}
\hline
\textbf{Work} & \textbf{Model} & \textbf{Result} &
\textbf{Limitation vs MakFleet} \\
\hline
Yu et al. \cite{yu2018} & STGCN & MAE 3.99 (METR-LA) &
Static graph; no semantic enrichment \\
\hline
Zhao et al. \cite{zhao2019} & T-GCN & MAE 2.37 (SZ-taxi) &
City-scale sensors; no campus micro-mobility \\
\hline
Newson \& Krumm \cite{newson2009} & HMM & 98.6\% road ID &
O(N$^2$) per step; no ontology integration \\
\hline
Dwork et al. \cite{dwork2006} & Diff. Privacy & Formal DP &
Aggregate only; no trajectory privacy \\
\hline
Pan et al. \cite{pan2024} & LLM+KG & Interoperability &
No real-time IoT pipeline or anomaly detection \\
\hline
\textbf{MakFleet} & GCN+LSTM & \textbf{AUC 0.9471 OOD} &
\textbf{Full semantic ST-KG with causal BI} \\
\hline
\end{tabular}
\end{table}


\section{System Architecture}

\subsection{Eight-Layer Locator Formalisation}

MakFleet is formalised as an \textbf{(L0)} operational safety and
spatio-temporal optimisation problem using \textbf{(L1)} weak
supervision from IoT telematics and OpenStreetMap to solve an
\textbf{(L2)} anomaly detection and demand forecasting task on
\textbf{(L3)} multivariate time-series and spatial graph data,
trained under an \textbf{(L4)} continual learning regime across
semester concept drift, under strict \textbf{(L5)}
privacy-preserving and out-of-distribution robustness constraints,
deployed via \textbf{(L6)} edge-to-cloud streaming into a
university-governed semantic warehouse, using \textbf{(L7)} ST-GNN
and rule-based symbolic models.

\subsection{Semantic Data Engineering Pipeline}

The pipeline operates in three continuous phases.

\textbf{Phase 1 -- IoT Ingestion and Provenance.} Bodaboda IoT
units stream GPS coordinates, 3-axis accelerometer readings, speed,
heading, and engine state at 5-second intervals. Every record
receives a SHA-256 cryptographic hash at ingestion. Batches are
chained into a Merkle tree whose root hash is stored in the
Knowledge Graph as a ProvenanceAnchor node, enabling post-ingestion
tamper detection.

\textbf{Phase 2 -- Semantic Map-Matching.} Raw GPS coordinates
carry up to 5~m Gaussian noise. A Hidden Markov Model
map-matcher---using a KD-tree spatial index over edge midpoints and
Viterbi decoding over a sliding window of 10 observations---snaps
each coordinate to the correct road edge. Mean map-matching
confidence achieved is 0.906.

\textbf{Phase 3 -- Event Enrichment.} Enriched records are
classified into six semantic event types using rule-based
thresholds: \texttt{speeding} (speed $> 30 \times 1.2$~km/h),
\texttt{harsh\_braking} ($a_x < -3.5$~m/s$^2$),
\texttt{rapid\_acceleration} ($a_x > 2.8$~m/s$^2$),
\texttt{idling}, \texttt{normal\_travel}, and
\texttt{safe\_stop}. The critical ontological distinction: a safe
stop near a mapped zebra crossing is loaded as a
\texttt{ContextualStop} node in Neo4j---a distinct label from
\texttt{TelematicsEvent}---enforcing the semantic separation at the
storage layer rather than at query time.

Before Knowledge Graph ingestion, GPS coordinates are subjected to
spatial k-anonymity ($k=5$) and Gaussian differential privacy noise
injection ($\varepsilon=1.0$). Role-Based Access Control provides
three permission tiers: admin, analyst, and public.

\subsection{Spatio-Temporal Knowledge Graph}

The ST-KG is stored in Neo4j and comprises five node labels
(Intersection, RoadSegment, Vehicle, Driver, TelematicsEvent/
ContextualStop) and five relationship types (ADJACENT\_TO,
LOCATED\_AT, CAUSED\_BY, DRIVES, PRECEDES). The campus road
network of 212 intersections and 476 road segments is retrieved
from OpenStreetMap via OSMnx \cite{boeing2017}. The full dataset
of 826,260 enriched records is publicly available at:
\texttt{kaggle.com/datasets/sebunyaronaldo/}\newline
\texttt{makfleet-bodaboda-telemetry-makerere-university}.

\subsection{ST-GNN Model Architecture}

The ST-GNN takes as input $T=4$ consecutive 15-minute graph
snapshots, each with node feature matrix
$\mathbf{X}_t \in \mathbb{R}^{N \times 6}$ (latitude, longitude,
betweenness centrality, degree centrality, event count, mean speed).

The spatial encoder applies two Graph Convolutional Network layers
per snapshot \cite{kipf2017}:
\begin{equation}
\mathbf{H}^{(l+1)} = \sigma\!\left(
  \tilde{\mathbf{D}}^{-\frac{1}{2}}
  \tilde{\mathbf{A}}\,
  \tilde{\mathbf{D}}^{-\frac{1}{2}}
  \mathbf{H}^{(l)}\mathbf{W}^{(l)}
\right)
\end{equation}
where $\tilde{\mathbf{A}} = \mathbf{A} + \mathbf{I}_N$ and
$\tilde{\mathbf{D}}$ is its degree matrix. The temporal aggregator
stacks the four snapshot embeddings and processes them through a
per-node LSTM:
\begin{equation}
\hat{p}_i = \sigma_s\!\left(
  \mathbf{W}_2\,\text{ReLU}(\mathbf{W}_1\,\mathbf{h}_{T,i}
  +\mathbf{b}_1)+\mathbf{b}_2\right)
\end{equation}
producing per-node anomaly probability $\hat{p}_i \in [0,1]$.
Weighted Binary Cross-Entropy with positive-class weight
$w^{+}=19.0$ addresses the 10:1 class imbalance. The model is
optimised with Adam ($\text{lr}=10^{-3}$) and Cosine Annealing.

Training strictly uses Semester~1 data; Semester~2 data serves as a
temporal out-of-distribution test, enforcing causal validity and
preventing data leakage.


\section{Results and Discussion}

\subsection{Exploratory Data Analysis}

EDA on the 826,260-record dataset reveals three key design insights.
Peak demand occurs at 08:00 aligned with morning lectures, with a
secondary peak at 17:00. Anomaly rate is elevated in early-morning
hours when low traffic density enables higher speeds. Speeding is
the most frequent anomaly class (91.2\% of records are normal
travel; 8.8\% anomalous). Driver risk profiles show strong positive
correlation with observed anomaly rates, validating the simulation.

\subsection{ST-GNN Evaluation}

Table~\ref{tab:results} reports performance across all evaluation
splits after 10 training epochs.

\begin{table}[t]
\caption{ST-GNN Evaluation Results}
\label{tab:results}
\centering
\begin{tabular}{|l|c|c|c|c|}
\hline
\textbf{Split} & \textbf{AUC} & \textbf{F1} &
\textbf{Prec.} & \textbf{Recall} \\
\hline
Validation (Sem.\,1) & 0.9268 & 0.9067 & 0.8293 & 1.000 \\
\hline
Temporal OOD (Sem.\,2) & \textbf{0.9471} & 0.9045 & 0.8256 & 1.000 \\
\hline
Spatial OOD (East) & ---$^{*}$ & 1.000 & 1.000 & 1.000 \\
\hline
\multicolumn{5}{l}{\footnotesize $^{*}$AUC undefined: all samples
positive class in eastern zone.}
\end{tabular}
\end{table}

The central finding is that the Temporal OOD AUC (0.9471) exceeds
the Validation AUC (0.9268), demonstrating that the GCN spatial
encoder learns generalisable road-topology risk patterns from
Semester~1 that transfer to unseen Semester~2 data. Recall of 1.000
across all splits confirms zero missed anomalies in every evaluation
condition---the operationally critical metric for campus safety.

\subsection{Causal Business Intelligence}

When the ST-GNN flags a driver, the system executes a 2-hop
ego-graph retrieval from Neo4j via ADJACENT\_TO traversal, rendering
the causal subgraph on a Folium campus map with colour-coded nodes
and generating a written explanation. For example: \textit{``Harsh
braking detected: deceleration $-4.2$~m/s$^2$ at 47~km/h. No zebra
crossing found within 20~m in campus ontology---classified as
reckless.''} The system is live at:
\texttt{makfleet.streamlit.app}.


\section{Conclusion}

MakFleet demonstrates that an ontology-driven Spatio-Temporal
Knowledge Graph can transform raw bodaboda IoT telemetry into
actionable, auditable campus safety evidence. The neuro-symbolic
architecture---combining GCN spatial encoding with rule-based
semantic event enrichment---achieves Temporal OOD AUC of 0.9471
with perfect recall, validating generalisation across semester
concept drift. Future work will deploy physical IoT units on actual
Makerere campus bodabodas and implement federated learning across
multiple Ugandan universities.


\section*{Acknowledgement}

The authors thank Dr.\ Ggaliwango Marvin for supervision and the
Department of Computer Science, COCIS, Makerere University for
providing the research environment.


\begin{thebibliography}{00}

\bibitem{kimball2013}
R. Kimball and M. Ross, \textit{The Data Warehouse Toolkit}, 3rd ed.
Hoboken, NJ: Wiley, 2013.

\bibitem{pelekis2014}
N. Pelekis and Y. Theodoridis, ``Mobility Data Management and
Exploration,'' \textit{Springer Briefs in Computer Science}, 2014.

\bibitem{zhu2022}
L. Zhu et al., ``Graph Neural Networks for Anomaly Detection in
Smart Transportation,'' \textit{IEEE Internet of Things Journal},
vol.~9, no.~23, 2022.

\bibitem{yu2018}
B. Yu, H. Yin, and Z. Zhu, ``Spatio-temporal graph convolutional
networks: A deep learning framework for traffic forecasting,''
in \textit{Proc. IJCAI}, 2018.

\bibitem{zhao2019}
L. Zhao et al., ``T-GCN: A temporal graph convolutional network for
traffic prediction,'' \textit{IEEE Trans. Intell. Transp. Syst.},
2019.

\bibitem{newson2009}
P. Newson and J. Krumm, ``Hidden Markov map matching through noise
and sparseness,'' in \textit{Proc. ACM SIGSPATIAL GIS}, 2009.

\bibitem{dwork2006}
C. Dwork et al., ``Calibrating noise to sensitivity in private data
analysis,'' in \textit{Proc. TCC}, 2006.

\bibitem{demontjoye2013}
Y. A. de Montjoye et al., ``Unique in the crowd: The privacy bounds
of human mobility,'' \textit{Scientific Reports}, vol.~3, 2013.

\bibitem{pan2024}
J. Pan et al., ``Unifying large language models and knowledge
graphs: A roadmap,'' \textit{IEEE TKDE}, 2024.

\bibitem{boeing2017}
G. Boeing, ``OSMnx: New methods for acquiring, constructing,
analyzing, and visualizing complex street networks,''
\textit{Computers, Environment and Urban Systems}, vol.~65,
pp.~126--139, 2017.

\bibitem{kipf2017}
T. N. Kipf and M. Welling, ``Semi-supervised classification with
graph convolutional networks,'' in \textit{Proc. ICLR}, 2017.

\end{thebibliography}

\end{document}
